\section{Discussion and Future Work}

\label{sec:discussion}

\subsection{Dynamic Mid-Session Tier Migration}

The current tier transition model upgrades but does not freely migrate
between arbitrary tiers mid-session. A session on Tier 3 cannot downgrade
to Tier 2 without risking state loss if the running process has no safe
checkpoint. Future work will develop a \textbf{live migration protocol}
analogous to VM live migration~\cite{clark2005live}: the agent's execution
state (open files, running processes, active connections) is serialized
incrementally while the target environment is prepared, minimizing migration
time and eliminating visible downtime.

\subsection{Multiplayer Agent Collaboration}

The current architecture supports multi-agent execution within a single
session but does not address multiple human users collaborating on the same
agent session in real time. We are developing a \textbf{collaborative session
model} inspired by CRDTs~\cite{shapiro2011conflict}: each participant's
instructions are modeled as operations on a shared session document, with
conflict resolution handled by the orchestrator agent. This enables scenarios
such as a developer and a product manager jointly directing an agent to
implement a feature, with each able to inspect the agent's reasoning and
interject corrections.

\subsection{Autonomous Operation with Configurable Trust Boundaries}

The current system requires human approval for certain actions (large cost
expenditures, irreversible system changes, external communications). Future
work will introduce a \textbf{trust boundary framework} that allows
organizations to configure the set of actions an agent may take autonomously
vs.\ those that require human approval, with cryptographically signed audit
logs for all autonomous actions. This is related to AI safety work on
corrigibility~\cite{soares2015corrigibility} and controllability, but
tailored to the practical constraints of production deployment.

\subsection{Federated Self-Improvement}

The current self-improvement loops operate on a single organization's data.
The Hanzo DSO (Decentralized Semantic Optimization) protocol~\cite{hanzo2024dso}
defines a privacy-preserving mechanism for sharing improvements across
organizations using zero-knowledge proofs and federated aggregation.
Integration of DSO with the UAH's Layer 6 would allow the harness to
improve from collective experience while preserving per-organization data
isolation---analogous to federated learning~\cite{mcmahan2017communication}
but operating on structured semantic priors rather than gradient updates.

\subsection{Limitations}

The UAH has several notable limitations that bound its applicability:

\begin{enumerate}[leftmargin=1.5em]
    \item \textbf{Task scope:} The harness performs best on tasks with
          measurable outcomes (code tests, metric targets). Open-ended creative
          tasks lack the feedback signal required for the self-improvement layer.
    \item \textbf{Latency:} Tier upgrades and agent handoffs introduce latency
          overhead (1--90 seconds) that is unacceptable for real-time
          interactive applications. Tier 0 (sub-second response) must be used
          for latency-critical workflows.
    \item \textbf{Cost predictability:} Because the harness dynamically selects
          compute tiers, costs per session can vary by up to $100\times$
          depending on task complexity. Budget caps and cost estimation help
          but do not eliminate cost uncertainty.
    \item \textbf{Model dependence:} Despite routing across 100+ providers,
          the overall system performance is bounded by the capability of the
          underlying frontier models. Architectural improvements cannot
          substitute for model capability on tasks requiring reasoning beyond
          model capacity.
\end{enumerate}

% ============================================================================
% 10. CONCLUSION
% ============================================================================
